\chapter{Phase 1 --- Algorithm--Hardware Co-Design and the Fixed-Point Golden Model}\label{ch:codesign}

\weeks{3--7} \quad Milestone M1

This is the most important phase of the project. You turn the floating-point
thesis code into a \emph{hardware specification}: a sample rate, frame size,
number format and function approximation for every signal, each justified by an
experiment. You also build a Python model that reproduces, bit for bit, what the
hardware will compute. Phase~2 then implements that model in RTL, and Phase~3
proves the RTL matches it.

\begin{whybox}
Phase~1 answers RQ1 (precision) and RQ2 (latency), tests H1 and H2, and gives a
first answer to RQ4 (H3, the quantised classifier head). It is also where most
of the \emph{research} in the project happens. RTL implements decisions; this
phase makes them. A hardware engineer who can say \enquote{16 bits here, because
this sweep shows that 14 costs 0.08 PESQ} is worth far more than one who picked
16 because it is a round number.
\end{whybox}

\begin{center}\small
\begin{tabularx}{\textwidth}{@{}LL@{}}
\toprule
\textbf{Inputs} & \textbf{Outputs} \\
\midrule
M0: metrics harness, float baseline, development and test sets; \cref{ch:workload}
and \cref{app:factsheet}
& \file{golden/} bit-accurate model; sweep results and plots; a specification
document; golden test vectors; NLP head quantisation results; repository tag
\code{spec-v1} \\
\bottomrule
\end{tabularx}
\end{center}

\begin{conceptbox}{Bit-accurate versus bit-exact}
A \emph{bit-accurate} model computes with exactly the integers the hardware
will hold. Every rounding, saturation and table lookup is explicit, and none of
it is left to numpy's floating point. When the RTL later produces the same
integers on the same inputs, the two are \emph{bit-exact}. That is what
verification checks (\cref{ch:verification}). Floating-point code can only ever
be \enquote{close}, and \enquote{close} gives you nothing to test against when
debugging hardware.
\end{conceptbox}

% =========================================================================
\section{Steps}\label{sec:codesign:steps}

\begin{stepbox}[label=st:codesign:profile]{Instrument the float model and profile dynamic range}
Before choosing any number format you need to know the range of every signal.
Subclass or wrap \code{IMCRA} and \code{ParametricWienerFilter} so each
intermediate quantity is recorded per frame, without changing any result:
\code{power}, $S$, $S_\text{min}$, $\gamma_\text{min}$, $\zeta$, the indicator,
$\tilde S$, $\tilde S_\text{min}$, $q$, $\gamma$, $\nu$, $p$, $\lambda_d$, $\xi$,
$G$, and the output samples.

Run the development set through it. For each signal, record the minimum, the
maximum, the 0.1st and 99.9th percentiles, and a histogram of $\log_2|x|$.

\textbf{Check:} with the probes enabled, the output is identical to the
uninstrumented model (assert this). You can fill the first two columns of the
specification table (\cref{tab:codesign:spec}). Expect power-like quantities to
span more than \qty{100}{\dB}. That is the first sign they want a logarithmic
or floating-point format rather than plain fixed point.

\textbf{Effort:} 1--2 days.
\end{stepbox}

\begin{stepbox}[label=st:codesign:frames]{Frame-configuration study (RQ2, H2)}
Sweep the configurations in \cref{tab:workload:configs}, plus the thesis
configuration with and without zero padding. The thesis \code{AudioConfig}
already accepts other sizes, provided \code{hop = window/2} and
\code{block = hop}. For \qty{16}{\kHz} runs, resample the input and the clean
reference with \code{resample\_poly}.

The step people forget is rescaling the time constants. IMCRA's smoothing
factors, the decision-directed factor and the minimum-search window are all
defined per frame. Keep their \emph{time constants} fixed when the frame period
$T=H/f_s$ changes to $T'$:

\begin{lstlisting}[style=py,caption={Keeping time constants fixed across frame configurations.}]
from dataclasses import replace

def rescale(imcra, wiener, T_old, T_new):
    r = T_new / T_old                      # ratio of frame periods
    a = lambda alpha: alpha ** r           # tau = -T/ln(alpha) stays fixed
    v_new = max(1, round(imcra.v / r))     # keep U = 8 sub-windows, D*T fixed
    return (replace(imcra, alpha_s=a(imcra.alpha_s), alpha_d=a(imcra.alpha_d), v=v_new),
            replace(wiener, dd_alpha=a(wiener.dd_alpha)))
\end{lstlisting}

For each configuration, record PESQ, STOI and SI-SDR on the development set,
together with the algorithmic latency and the state memory. Plot quality against
latency and mark the configurations that dominate the others: they form the
Pareto front (\cref{fig:codesign:schematic}, right).

\textbf{Check:} the thesis configuration reproduces the baseline exactly. You
can state H2's outcome: what removing the zero padding costs in PESQ. You have
chosen a design point and written down why (decision point, week~4,
\cref{ch:plan}).

\textbf{Effort:} 3--4 days.
\end{stepbox}

\begin{stepbox}[label=st:codesign:fxlib]{Write a bit-true fixed-point library}
Every fixed-point operation in the model goes through a few helper functions.
The RTL later implements exactly the same rounding and saturation. Use numpy
\code{int64} arrays, which have enough headroom for the products of two 32-bit
values.

\begin{lstlisting}[style=py,caption={The heart of \texttt{golden/fixedpoint.py}.}]
import numpy as np

def limits(word, signed):
    return (-(1 << (word - 1)), (1 << (word - 1)) - 1) if signed else (0, (1 << word) - 1)

def saturate(q, word, signed=True):
    lo, hi = limits(word, signed)
    return np.clip(q, lo, hi).astype(np.int64)

def quantize(x, word, frac, signed=True):
    """Real value(s) -> integer(s) in Q(word-frac).frac, round half up, saturate."""
    q = np.floor(np.asarray(x, dtype=np.float64) * (1 << frac) + 0.5)
    return saturate(q, word, signed)

def requant(q, frac_in, word_out, frac_out, signed=True):
    """Change the binary point of integer q, rounding half up, then saturate."""
    shift = frac_in - frac_out
    if shift > 0:
        q = (q + (1 << (shift - 1))) >> shift      # arithmetic shift = floor
    else:
        q = q << (-shift)
    return saturate(q, word_out, signed)

def mul(a, fa, b, fb, word_out, frac_out, signed=True):
    """Exact integer product, then one rounding and one saturation."""
    return requant(a.astype(np.int64) * b.astype(np.int64), fa + fb,
                   word_out, frac_out, signed)

def to_real(q, frac):
    return q.astype(np.float64) / (1 << frac)
\end{lstlisting}

Choose \emph{one} rounding mode and use it everywhere. Round-half-up (add half
an LSB, then shift right) costs one adder in hardware, and that is what the code
above does.

\textbf{Check:} unit tests for exact representable values, rounding at the
half-LSB boundary for positive and negative numbers, saturation at both limits,
and \code{mul} against exact rational arithmetic (Python \code{fractions}) on
random inputs. \textbf{Effort:} 1 day.
\end{stepbox}

\begin{stepbox}[label=st:codesign:blocks]{Build the bit-accurate model, block by block, in hardware order}
Write one module per hardware block: \file{golden/fx_stft.py},
\file{golden/fx_imcra.py}, \file{golden/fx_wiener.py}, and
\file{golden/fx_pipeline.py} to connect them. Quantise one block at a time and
compare against float after each change. If you quantise everything at once and
the quality collapses, you will not know why.

\begin{enumerate}
  \item \textbf{Window and framing.} The $\sqrt{\text{Hann}}$ window becomes a
  ROM of quantised coefficients. By symmetry only half of it needs storing.
  \item \textbf{FFT.} Model the \emph{hardware} algorithm exactly, not
  \code{np.fft}. Write a radix-2 decimation-in-time FFT~\cite{cooley1965fft}
  with quantised twiddle factors, and round after every butterfly. Add
  \emph{block floating point} (O8): after each stage, if any value has grown
  into the guard bits, shift the whole block right by one and increment a shared
  exponent. The inverse FFT reuses the forward FFT through conjugation:
  $\mathrm{IFFT}(X)=\overline{\mathrm{FFT}(\overline{X})}/N$.
  \item \textbf{Power.} $|X_k|^2$ doubles the dynamic range in bits. This is
  where you move to a logarithmic or small floating-point format for power-like
  quantities (\cref{st:codesign:log}).
  \item \textbf{IMCRA and Wiener.} Implement O3 (cache $\xi/(1+\xi)$), O4
  (shared reciprocals via seed table + Newton--Raphson) and O1 (gain LUT,
  \cref{st:codesign:lut}).
  \item \textbf{Tilt and output gain} (O2), the gain multiply, the inverse FFT,
  the synthesis window and overlap-add.
\end{enumerate}

\textbf{Check:} with very wide formats (for example 32 fractional bits
everywhere and large tables), the model matches float to within about
0.01 PESQ and 0.001 STOI. That shows the \emph{structure} is
right before you start narrowing word lengths.

\textbf{Effort:} 1.5--2 weeks. The FFT and block floating point take the most
time.
\end{stepbox}

\begin{stepbox}[label=st:codesign:log]{Decide the representation of power-like quantities (O5)}
$|X|^2$, $S$, $\tilde S$, their minima and $\lambda_d$ span a huge range.
There are three candidate representations:
\begin{enumerate}
  \item \textbf{Wide fixed point} (about 32 bits or more). Simple, but expensive
        in memory, and memory is 23 words per bin.
  \item \textbf{Small custom floating point}, for example a 10-bit mantissa with
        a 6-bit exponent. Its relative precision is constant across the range.
  \item \textbf{Logarithmic}, for example $\log_2$ as a signed Q6.10 value.
        Multiplications and divisions become additions and subtractions, and
        comparisons stay exact.
\end{enumerate}

The minimum trackers, which hold 18 of the 23 state words, are a natural fit for
the log domain. A minimum is unchanged by any monotone map, so
$\min(\log a,\log b)=\log\min(a,b)$ exactly. The threshold tests in
\file{imcra.py} also survive unchanged: $\gamma_\text{min}<\gamma_0$ becomes
$\log_2 P - \log_2(B_{\min}S_\text{min}) < \log_2\gamma_0$.

\emph{Recursive smoothing does not survive.} $\alpha\log S+(1-\alpha)\log P$ is
the log of a \emph{geometric} mean, which is biased low compared with the
arithmetic mean that IMCRA expects. For a single exponentially distributed
periodogram bin the bias is about \qty{2.5}{\dB}, the Euler--Mascheroni
constant in nepers. Smoothing reduces it, but it does not vanish.

Evaluate two options:
\begin{enumerate}[label=(\alph*)]
  \item keep the smoothed quantities $S$, $\tilde S$, $\lambda_d$ in a small
        floating-point format, and store only the tracker minima as
        low-precision log values;
  \item smooth in the log domain and re-fit the bias factor $B_{\min}$ so the
        noise estimate is unbiased on the development set.
\end{enumerate}

\textbf{Check:} for your chosen option, the noise PSD estimate on the
development set has a median error under about \qty{0.5}{\dB} against float,
and the end metrics meet RQ1. \textbf{Effort:} 3--4 days.
\end{stepbox}

\begin{stepbox}[label=st:codesign:lut]{Gain LUT and control-plane mapping (O1, O2)}
Replace $(\xi/(1+\xi))^\beta$ by a table that the control plane computes
whenever $\beta$ or the gain floor changes:

\begin{lstlisting}[style=py,caption={Gain table generation (control plane) and lookup (data plane).}]
import numpy as np
from golden.fixedpoint import quantize

LOG2_STEP = 1 / 8            # LUT spacing: 1/8 of a log2 unit (about 0.38 dB)
LOG2_MIN = -9.0              # about -27 dB, just below xi_min = -25 dB
N_LUT = 256                  # covers 32 log2 units, about 96 dB

def make_gain_lut(beta, floor):
    log2_xi = LOG2_MIN + LOG2_STEP * np.arange(N_LUT)
    xi = 2.0 ** log2_xi
    g = np.maximum((xi / (1 + xi)) ** beta, floor)
    return quantize(g, word=16, frac=13, signed=False)     # UQ3.13

def lut_index(log2_xi_q, frac=10):
    # log2(xi) held as signed Q6.10. A power-of-two step makes the index a
    # subtraction and a shift -- no multiplier.
    idx = (log2_xi_q - int(LOG2_MIN * (1 << frac))) >> (frac - 3)
    return np.clip(idx, 0, N_LUT - 1)
\end{lstlisting}

Choosing the table spacing as a power of two in $\log_2$ units is a small
co-design decision with a real payoff: the table index is a bit slice of the
log-domain $\xi$, so no multiplier is needed. For the tilt (O2), precompute
$r=10^{\text{tilt}/(20(K-1))}$ and form $g_k=g_\text{out}\,r^k$ with one multiply
per bin. Repeated multiplication accumulates rounding error over the $K$ bins,
so check the error at the last bin. Alternatively, add a constant per bin in the
log domain.

Sweep the table size (64--1024 entries), the $\xi$ range, and linear
interpolation versus none.

\textbf{Check:} applying the thesis's own parameter changes (for example
$\beta$ from 1.0 to 1.94, then a tilt step) produces the same audio as the float
model to within the RQ1 target. The table swap at a frame boundary produces no
click: compare the difference signal around the switch.
\textbf{Effort:} 2--3 days.
\end{stepbox}

\begin{stepbox}[label=st:codesign:sweeps]{Word-length and approximation sweeps (RQ1, H1)}
Now narrow the formats. Use the development set, and a fixed subset of it while
iterating.
\begin{enumerate}
  \item \textbf{One factor at a time.} Vary one signal's word length (or one
  table's size), keeping all others wide. Find the knee where quality starts to
  fall (\cref{fig:codesign:schematic}, left).
  \item \textbf{Joint.} Set every signal a little above its knee and run the
  whole model. Quantisation errors add up, so you will usually need a bit or two
  more than the one-at-a-time results suggest.
  \item \textbf{Nonlinear functions.} Size the reciprocal seed table and the
  number of Newton--Raphson iterations. One iteration roughly doubles the number
  of correct bits. Size the $e^{-\nu}$ table (O7) over $\nu\in[0,16]$, or the
  sigmoid rewrite of \cref{eq:workload:sigmoid}. Derive an accuracy target for
  each function from the end metric, not from intuition.
\end{enumerate}

\begin{center}
\begin{tikzpicture}
\begin{groupplot}[group style={group size=2 by 1,horizontal sep=18mm},
  width=0.47\textwidth,height=42mm,
  tick label style={font=\sffamily\scriptsize},label style={font=\sffamily\scriptsize},
  axis lines=left,xtick=\empty,ytick=\empty,clip=false]
\nextgroupplot[xlabel={word length (bits)},ylabel={quality loss vs float},xmin=6,xmax=24,ymin=0,ymax=1.1]
  \addplot[brand,thick,domain=6:24,samples=60] {exp(-(x-6)*0.55)};
  \draw[pitc,dashed] (axis cs:6,0.1) -- (axis cs:24,0.1) node[right,font=\sffamily\tiny,text=pitc]{target};
  \draw[-{Stealth},pitc] (axis cs:13,0.55) node[above,font=\sffamily\tiny,text=pitc]{knee} -- (axis cs:10.3,0.12);
\nextgroupplot[xlabel={latency},ylabel={quality},xmin=0,xmax=10,ymin=0,ymax=10]
  \addplot[only marks,mark=*,mark size=1.6pt,conceptc] coordinates {(2,3.5)(3.2,4.4)(4,6.3)(5.5,5.6)(6,7.5)(7.3,7.3)(8.2,8.4)(9,7.9)};
  \addplot[brand,thick] coordinates {(2,3.5)(3.2,4.4)(4,6.3)(6,7.5)(8.2,8.4)};
  \node[font=\sffamily\tiny,text=brand,anchor=west] at (axis cs:6.6,6.3) {Pareto front};
\end{groupplot}
\end{tikzpicture}
\captionof{figure}{Schematic --- illustrative shape, not data. Left: quality loss
against word length typically falls steeply and then flattens. Pick the smallest
word length past the knee that meets the target. Right: configurations that are
not beaten on both latency and quality form the Pareto front.}
\label{fig:codesign:schematic}
\end{center}

\textbf{Check:} a joint configuration meets RQ1 on the development set. A table
lists every signal with its format and the experiment that justifies it. You can
state H1's outcome. \textbf{Effort:} 1 week (mostly computer time; run the
sweeps overnight).
\end{stepbox}

\begin{stepbox}[label=st:codesign:nlp]{Quantise the NLP heads (H3, first part of RQ4)}
Use the thesis pipeline's cached sentence embeddings (\code{SentenceEmbedder}
writes them to its \code{cache\_dir}) and the thesis's fixed train/test split.
Fit logistic regression and the thesis's linear SVM on the same split, then
quantise the weights to int8 with one scale per row:

\begin{lstlisting}[style=py,caption={Post-training int8 quantisation of a linear head.}]
import numpy as np

def quantize_head(W, b, X):
    s_w = np.abs(W).max(axis=1, keepdims=True) / 127        # per-row scale
    Wq = np.clip(np.round(W / s_w), -127, 127).astype(np.int8)
    s_x = 1.0 / 127                                         # |x_i| <= 1 (L2-normalised)
    Xq = np.clip(np.round(X / s_x), -127, 127).astype(np.int8)
    acc = Xq.astype(np.int32) @ Wq.astype(np.int32).T      # int32 accumulators
    return acc * (s_w.T * s_x) + b                          # dequantised scores

# Logistic regression: 6 rows -> argmax.  Linear SVC (one-vs-one): 15 rows,
# then the same pairwise voting as sklearn -- check your voting code on the
# float weights first: it must reproduce clf.predict exactly.
\end{lstlisting}

Report macro-F1 (float and int8), the multiply--accumulate count (\num{4608}
against \num{11520}), and the severity MAE with the Ridge weights quantised the
same way. Replace the SVM's Platt/pairwise-coupling probability with a
margin-based confidence (O9). For example, use the difference between the top
two scores, and pick its threshold on the training folds so the confidence gate
behaves like the thesis's.

\textbf{Check:} the float numbers reproduce the thesis (test macro-F1 0.957 for
the SVM). You can state H3's outcome. \textbf{Effort:} 2 days.
\end{stepbox}

\begin{stepbox}[label=st:codesign:freeze]{Freeze the specification and export golden vectors}
Write \file{docs/spec.md}: the chosen frame configuration, every signal's format
(\cref{tab:codesign:spec}), every table's size and contents rule, the rounding
mode, the block-floating-point rule, and the parameter-update semantics
(table bank swap at a frame boundary).

Then export \emph{golden vectors} with \file{golden/vectors.py} for the
verification phase. For each test signal (the synthetic set plus a few
VoiceBank utterances), write one file per probe point: the block inputs and
outputs per frame, as integers, in \code{.npy} or one-hex-word-per-line
\code{.hex} files. These are the contract between Python and RTL.

\begin{center}\footnotesize
\captionof{table}{Specification table template --- fill one row per signal and table.}
\label{tab:codesign:spec}
\begin{tabularx}{\textwidth}{@{}llllL@{}}
\toprule
Signal & Symbol & Format & Measured range & Evidence (experiment) \\
\midrule
Input sample & $x[n]$ & Q1.15 & & \file{experiments/...} \\
FFT data path & $X_k$ & & & \\
Power & $|X_k|^2$ & & & \\
Tracker minima & $S_\text{min}$, $\tilde S_\text{min}$ & & & \\
Noise PSD & $\lambda_d$ & & & \\
A priori SNR & $\xi$ & & & \\
Gain & $G_k$ & UQ3.13 & & \\
Gain LUT & --- & 256 $\times$ 16 bit & & \\
\ldots & & & & \\
\bottomrule
\end{tabularx}
\end{center}

\textbf{Check:} a fresh clone can regenerate every vector and every plot from
one command. Tag the repository \code{spec-v1}. \textbf{Effort:} 2 days.
\end{stepbox}

% =========================================================================
\section{Pitfalls}\label{sec:codesign:pitfalls}

\begin{pitfallbox}[Pitfalls in Phase 1]
\begin{itemize}
  \item \textbf{Using \code{np.fft} in the \enquote{bit-accurate} model.} The
        hardware FFT rounds after every stage. If the model does not, the RTL
        can never match it.
  \item \textbf{Silent float leaks.} A single \code{np.log2} or \code{/} on a
        fixed-point path makes the model no longer bit-accurate. Grep for them,
        and assert integer dtypes at block boundaries.
  \item \textbf{Comparing frame configurations without rescaling time
        constants.} You would be comparing different algorithms.
  \item \textbf{Optimising one signal at a time and stopping there.} Errors add
        up. Always finish with a joint run.
  \item \textbf{Tuning on the test set.} Use the development set for every
        choice.
  \item \textbf{Scope creep.} Phase~1 is fascinating and could last forever.
        It ends in week~7 (\cref{ch:plan}).
\end{itemize}
\end{pitfallbox}

\begin{readbox}
Fixed-point arithmetic and quantisation noise:
Widrow and Koll\'ar~\cite{widrow2008quantization}. Word-length optimisation:
Constantinides \etal~\cite{ch12-13:constantinides2003wordlength}. Elementary
function approximation (tables, polynomials, Newton--Raphson):
Muller~\cite{muller2016elementary}. IMCRA itself:
Cohen~\cite{cohen2003imcra}. Integer-only inference:
Jacob \etal~\cite{jacob2018quantization}.
\end{readbox}

\begin{donebox}[Milestone M1 is reached when]
\begin{checklist}
  \item The dynamic-range profile and the frame-configuration study are complete and plotted.
  \item The time-constant rescaling is implemented and tested.
  \item \file{golden/} is bit-accurate (integer-only on fixed-point paths) and unit-tested.
  \item The chosen formats meet RQ1 on the development set; H1 and H2 have recorded outcomes.
  \item The gain LUT, tilt mapping and table-swap semantics are specified and tested.
  \item The NLP heads are quantised; H3 has a recorded outcome.
  \item \file{docs/spec.md} is complete, golden vectors are exported, and the repository is tagged \code{spec-v1}.
\end{checklist}
\end{donebox}
